\subsection{Expresiones Regulares (módulo re)}
\label{alg:2-2}

El módulo \texttt{re} permite buscar, validar, reemplazar y extraer información de cadenas mediante patrones.

En programación competitiva suele utilizarse cuando la entrada posee un formato irregular, contiene separadores complejos o es necesario validar cadenas.

\begin{lstlisting}[language=Python]
import re
\end{lstlisting}

\vspace{0.3cm}

\subsubsection{Funciones principales}

\begin{center}
\begin{tabular}{|l|p{5cm}|}
\hline
Función & Descripción \\ \hline
re.match() & Busca únicamente al inicio del string. \\ \hline
re.search() & Busca la 1ra coincidencia en cualquier posición. \\ \hline
re.findall() & Devuelve todas las coincidencias en una lista. \\ \hline
re.finditer() & Devuelve un iter con todos los matches. \\ \hline
re.fullmatch() & Verifica que toda la cadena cumpla patrón. \\ \hline
re.sub() & Reemplaza coincidencias. \\ \hline
re.split() & Divide una cadena utilizando un patrón. \\ \hline
re.compile() & Compila una regex para reutilizarla. \\ \hline
\end{tabular}
\end{center}

\vspace{0.3cm}

\subsubsection{Símbolos más utilizados}

\begin{center}
\begin{tabular}{|c|l|}
\hline
Patrón & Significado \\ \hline
. & Cualquier carácter excepto salto de línea \\ \hline
\regex{\textbackslash d} & Dígito \\ \hline
\regex{\textbackslash D} & No dígito \\ \hline
\regex{\textbackslash w} & Letra, número o \_ \\ \hline
\regex{\textbackslash W} & No alfanumérico \\ \hline
\regex{\textbackslash s} & Espacio en blanco \\ \hline
\regex{\textbackslash S} & No espacio \\ \hline
+ & Una o más repeticiones \\ \hline
* & Cero o más repeticiones \\ \hline
? & Cero o una repetición \\ \hline
\{n\} & Exactamente n repeticiones \\ \hline
\{n,m\} & Entre n y m repeticiones \\ \hline
[] & Conjunto de caracteres \\ \hline
[\^{}] & Negación del conjunto \\ \hline
() & Grupo de captura \\ \hline
(?:) & Grupo sin captura \\ \hline
| & OR lógico \\ \hline
\^{} & Inicio del string \\ \hline
\$ & Final del string \\ \hline
\end{tabular}
\end{center}

\vspace{0.3cm}

\subsubsection{Ejemplos básicos}

\begin{lstlisting}[language=Python]
import re
nums = re.findall(r'-?\d+', s)
nums = list(map(int, nums))
tokens = re.split(r'\s+', s.strip())
m = re.match(r'^(\d+)-(\d+)$', s)
if m:
    a = int(m.group(1))
    b = int(m.group(2))
limpio = re.sub(r'[^a-zA-Z0-9]', '', s)
\end{lstlisting}

\vspace{0.3cm}

\subsubsection{Patrones útiles para ICPC}

\begin{center}
\begin{tabular}{|l|l|}
\hline
Patrón & Detecta \\ \hline
\texttt{\textbackslash d+} & Enteros positivos \\ \hline
\texttt{-?\textbackslash d+} & Enteros con signo \\ \hline
\texttt{\textbackslash d+\textbackslash.\textbackslash d+} & Números decimales \\ \hline
\texttt{[A-Za-z]+} & Palabras \\ \hline
\texttt{[A-Za-z0-9]+} & Alfanuméricos \\ \hline
\texttt{[aeiou]} & Vocales \\ \hline
\regex{[\string^aeiou]} & Todo excepto vocales \\ \hline
\texttt{\textbackslash s+} & Uno o más espacios \\ \hline
\texttt{\textbackslash w+} & Identificadores \\ \hline
\texttt{.+} & Cualquier texto \\ \hline
\end{tabular}
\end{center}

\vspace{0.3cm}

\subsubsection{Combinaciones frecuentes}

Extraer todos los enteros.

\begin{lstlisting}[language=Python]
re.findall(r'-?\d+', s)
\end{lstlisting}

Extraer únicamente números positivos.

\begin{lstlisting}[language=Python]
re.findall(r'\d+', s)
\end{lstlisting}

Extraer palabras.

\begin{lstlisting}[language=Python]
re.findall(r'[A-Za-z]+', s)
\end{lstlisting}

Extraer palabras y números.

\begin{lstlisting}[language=Python]
re.findall(r'\w+', s)
\end{lstlisting}

Separar por múltiples espacios.

\begin{lstlisting}[language=Python]
re.split(r'\s+', s.strip())
\end{lstlisting}

Eliminar caracteres especiales.

\begin{lstlisting}[language=Python]
re.sub(r'[^A-Za-z0-9]', '', s)
\end{lstlisting}

Eliminar espacios repetidos.

\begin{lstlisting}[language=Python]
re.sub(r'\s+', ' ', s)
\end{lstlisting}

Validar un correo electrónico sencillo.

\begin{lstlisting}[language=Python]
re.fullmatch(r'\w+@\w+\.\w+', s)
\end{lstlisting}

Validar una fecha DD/MM/AAAA.

\begin{lstlisting}[language=Python]
re.fullmatch(r'\d{2}/\d{2}/\d{4}', s)
\end{lstlisting}

Validar un entero.

\begin{lstlisting}[language=Python]
re.fullmatch(r'-?\d+', s)
\end{lstlisting}

Buscar únicamente al inicio.

\begin{lstlisting}[language=Python]
re.match(r'ICPC', s)
\end{lstlisting}

Buscar en cualquier posición.

\begin{lstlisting}[language=Python]
re.search(r'ICPC', s)
\end{lstlisting}

\vspace{0.3cm}

\subsubsection{Compilar una expresión regular}

Cuando un mismo patrón será utilizado muchas veces, es recomendable compilarlo una única vez.
Esto evita recompilar la expresión regular en cada llamada.

\begin{lstlisting}[language=Python]
pat = re.compile(r'-?\d+')
for linea in texto:
    nums = pat.findall(linea)
\end{lstlisting}

\vspace{0.3cm}

\subsubsection{Complejidad}

La mayoría de operaciones tienen complejidad aproximadamente \bigO{n}
respecto al tamaño del texto, aunque algunos patrones muy complejos pueden provocar retroceso (\textit{backtracking}) y aumentar considerablemente el tiempo de ejecución.

En problemas de ICPC normalmente se utilizan patrones simples, por lo que el rendimiento suele ser lineal.

\vspace{0.3cm}

\subsubsection*{Resumen}

\begin{center}
\begin{tabular}{|l|l|}
\hline
Función & Uso principal \\ \hline
match & Buscar al inicio \\ \hline
search & Buscar primera coincidencia \\ \hline
findall & Obtener todas las coincidencias \\ \hline
finditer & Iterar sobre coincidencias \\ \hline
fullmatch & Validar una cadena completa \\ \hline
split & Separar mediante un patrón \\ \hline
sub & Reemplazar texto \\ \hline
compile & Reutilizar un patrón \\ \hline
\end{tabular}
\end{center}
